
\documentclass{article}
\usepackage{colortbl}
\usepackage{makecell}
\usepackage{multirow}
\usepackage{supertabular}

\begin{document}

\newcounter{utterance}

\centering \large Interaction Transcript for game `dond', experiment `semi\_en', episode 10 with Qwen3{-}32B{-}t0.0.
\vspace{24pt}

{ \footnotesize  \setcounter{utterance}{1}
\setlength{\tabcolsep}{0pt}
\begin{supertabular}{c@{$\;$}|p{.15\linewidth}@{}p{.15\linewidth}p{.15\linewidth}p{.15\linewidth}p{.15\linewidth}p{.15\linewidth}}
    \# & $\;$A & \multicolumn{4}{c}{Game Master} & $\;\:$B\\
    \hline

    \theutterance \stepcounter{utterance}  
    & & \multicolumn{4}{p{0.6\linewidth}}{
        \cellcolor[rgb]{0.9,0.9,0.9}{
            \makecell[{{p{\linewidth}}}]{
                \texttt{\tiny{[P1$\langle$GM]}}
                \texttt{You are playing a negotiation game in which you have to agree on how to divide a set of items among you and another player.} \\
\\ 
\texttt{Rules:} \\
\texttt{(a) You and the other player are given a set of items. Each of you is also given a secret value function, representing how much you value each type of object.} \\
\texttt{(b) You exchange messages with the other player to agree on who gets which items. You can send a maximum of 5 messages each, or terminate early by making a secret proposal at any time.} \\
\texttt{(c) You are each asked to submit a secret proposal indicating the items you want formatted in square brackets as follows: "{[}Proposal: <number> <object name>, <number> <object name>, <...>{]}"} \\
\texttt{(d) If your proposals are complementary, i.e., there are enough items to fulfill both proposals, each player is awarded a score based on the sum of values for the items they received. Otherwise, both of you get zero points.} \\
\texttt{(e) Your goal is to maximize the score you receive.} \\
\\ 
\texttt{Let us start.} \\
\\ 
\texttt{The set of available items is:} \\
\\ 
\texttt{1 bucket, 3 knives, 1 battery} \\
\\ 
\texttt{Your secret value function is:} \\
\\ 
\texttt{bucket: 2, knife: 2, battery: 2} \\
\\ 
\texttt{IMPORTANT: Your messages, unless it is the secret proposal, are directly transmitted to the other player, so do not include any response to the rules or text announcing your message. To make a secret proposal, use the indicated format. Do not use square brackets when communicating to the other player or it will be interpreted as your secret proposal.} \\
            }
        }
    }
    & \\ \\

    \theutterance \stepcounter{utterance}  
    & \multicolumn{4}{p{0.6\linewidth}}{
        \cellcolor[rgb]{0.9,0.9,0.9}{
            \makecell[{{p{\linewidth}}}]{
                \texttt{\tiny{[P1$\rangle$GM]}}
                \texttt{<think>} \\
\texttt{Okay, let's see. I need to negotiate the division of 1 bucket, 3 knives, and 1 battery with another player. My value function is 2 for each item, so I value all items equally. The other player's values are secret, so I don't know what they prefer. The goal is to maximize my score, which is the sum of the values of the items I get. If our proposals are complementary, we both get our scores; otherwise, zero.} \\
\\ 
\texttt{First, I should consider possible strategies. Since I value all items equally, maybe I should aim for a fair split. But since the other player's preferences are unknown, I need to figure out what they might want. The items are 1 bucket, 3 knives, 1 battery. Let's think about possible splits.} \\
\\ 
\texttt{If I take the bucket and battery, that's two items, leaving three knives. But the other player might want some of those. Alternatively, if I take some knives, maybe the other player can take the rest. But since there are three knives, maybe splitting them 2{-}1 could work. But how to communicate this?} \\
\\ 
\texttt{Since I can send up to 5 messages, maybe I should start by asking the other player what they value. But I can't ask directly because the messages are exchanged. Wait, the messages are directly transmitted, so I can ask them to propose their preferences. But they might not respond honestly. Alternatively, I can make a proposal that allows for a complementary split.} \\
\\ 
\texttt{Given that my value is 2 for each item, the maximum I can get is 2* (number of items I get). Since there are 5 items total (1+3+1), if I get 3 items, my score is 6. But the other player would need to get 2 items. However, the other player's value function might not allow that. But since the proposals need to be complementary, I need to make sure that the items I want don't overlap with what they want.} \\
\\ 
\texttt{Wait, the secret proposal is when we both submit our desired items. If our proposals are complementary, we get our items. So maybe I should propose taking items that the other player might not want. But how to know?} \\
\\ 
\texttt{Alternatively, since I value all items equally, maybe I should aim for a split where I get as many items as possible without conflicting. For example, if I take the bucket and battery, that's} \\
            }
        }
    }
    & & \\ \\

    \theutterance \stepcounter{utterance}  
    & & & \multicolumn{2}{p{0.3\linewidth}}{
        \cellcolor[rgb]{0.9,0.9,0.9}{
            \makecell[{{p{\linewidth}}}]{
                \texttt{\tiny{[GM$|$GM]}}
                \texttt{continue} \\
            }
        }
    }
    & & \\ \\

    \theutterance \stepcounter{utterance}  
    & & & \multicolumn{2}{p{0.3\linewidth}}{
        \cellcolor[rgb]{0.9,0.9,0.9}{
            \makecell[{{p{\linewidth}}}]{
                \texttt{\tiny{[GM$|$GM]}}
                \texttt{<think>} \\
\texttt{Okay, let's see. I need to negotiate the division of 1 bucket, 3 knives, and 1 battery with another player. My value function is 2 for each item, so I value all items equally. The other player's values are secret, so I don't know what they prefer. The goal is to maximize my score, which is the sum of the values of the items I get. If our proposals are complementary, we both get our scores; otherwise, zero.} \\
\\ 
\texttt{First, I should consider possible strategies. Since I value all items equally, maybe I should aim for a fair split. But since the other player's preferences are unknown, I need to figure out what they might want. The items are 1 bucket, 3 knives, 1 battery. Let's think about possible splits.} \\
\\ 
\texttt{If I take the bucket and battery, that's two items, leaving three knives. But the other player might want some of those. Alternatively, if I take some knives, maybe the other player can take the rest. But since there are three knives, maybe splitting them 2{-}1 could work. But how to communicate this?} \\
\\ 
\texttt{Since I can send up to 5 messages, maybe I should start by asking the other player what they value. But I can't ask directly because the messages are exchanged. Wait, the messages are directly transmitted, so I can ask them to propose their preferences. But they might not respond honestly. Alternatively, I can make a proposal that allows for a complementary split.} \\
\\ 
\texttt{Given that my value is 2 for each item, the maximum I can get is 2* (number of items I get). Since there are 5 items total (1+3+1), if I get 3 items, my score is 6. But the other player would need to get 2 items. However, the other player's value function might not allow that. But since the proposals need to be complementary, I need to make sure that the items I want don't overlap with what they want.} \\
\\ 
\texttt{Wait, the secret proposal is when we both submit our desired items. If our proposals are complementary, we get our items. So maybe I should propose taking items that the other player might not want. But how to know?} \\
\\ 
\texttt{Alternatively, since I value all items equally, maybe I should aim for a split where I get as many items as possible without conflicting. For example, if I take the bucket and battery, that's} \\
            }
        }
    }
    & & \\ \\

    \theutterance \stepcounter{utterance}  
    & & \multicolumn{4}{p{0.6\linewidth}}{
        \cellcolor[rgb]{0.9,0.9,0.9}{
            \makecell[{{p{\linewidth}}}]{
                \texttt{\tiny{[GM$\rangle$P2]}}
                \texttt{You are playing a negotiation game in which you have to agree on how to divide a set of items among you and another player.} \\
\\ 
\texttt{Rules:} \\
\texttt{(a) You and the other player are given a set of items. Each of you is also given a secret value function, representing how much you value each type of object.} \\
\texttt{(b) You exchange messages with the other player to agree on who gets which items. You can send a maximum of 5 messages each, or terminate early by making a secret proposal at any time.} \\
\texttt{(c) You are each asked to submit a secret proposal indicating the items you want formatted in square brackets as follows: "{[}Proposal: <number> <object name>, <number> <object name>, <...>{]}"} \\
\texttt{(d) If your proposals are complementary, i.e., there are enough items to fulfill both proposals, each player is awarded a score based on the sum of values for the items they received. Otherwise, both of you get zero points.} \\
\texttt{(e) Your goal is to maximize the score you receive.} \\
\\ 
\texttt{Let us start.} \\
\\ 
\texttt{The set of available items is:} \\
\\ 
\texttt{1 bucket, 3 knives, 1 battery} \\
\\ 
\texttt{Your secret value function is:} \\
\\ 
\texttt{bucket: 3, knife: 2, battery: 1} \\
\\ 
\texttt{IMPORTANT: Your messages, unless it is the secret proposal, are directly transmitted to the other player, so do not include any response to the rules or text announcing your message. To make a secret proposal, use the indicated format. Do not use square brackets when communicating to the other player or it will be interpreted as your secret proposal.} \\
\\ 
\\ 
\texttt{<think>} \\
\texttt{Okay, let's see. I need to negotiate the division of 1 bucket, 3 knives, and 1 battery with another player. My value function is 2 for each item, so I value all items equally. The other player's values are secret, so I don't know what they prefer. The goal is to maximize my score, which is the sum of the values of the items I get. If our proposals are complementary, we both get our scores; otherwise, zero.} \\
\\ 
\texttt{First, I should consider possible strategies. Since I value all items equally, maybe I should aim for a fair split. But since the other player's preferences are unknown, I need to figure out what they might want. The items are 1 bucket, 3 knives, 1 battery. Let's think about possible splits.} \\
\\ 
\texttt{If I take the bucket and battery, that's two items, leaving three knives. But the other player might want some of those. Alternatively, if I take some knives, maybe the other player can take the rest. But since there are three knives, maybe splitting them 2{-}1 could work. But how to communicate this?} \\
\\ 
\texttt{Since I can send up to 5 messages, maybe I should start by asking the other player what they value. But I can't ask directly because the messages are exchanged. Wait, the messages are directly transmitted, so I can ask them to propose their preferences. But they might not respond honestly. Alternatively, I can make a proposal that allows for a complementary split.} \\
\\ 
\texttt{Given that my value is 2 for each item, the maximum I can get is 2* (number of items I get). Since there are 5 items total (1+3+1), if I get 3 items, my score is 6. But the other player would need to get 2 items. However, the other player's value function might not allow that. But since the proposals need to be complementary, I need to make sure that the items I want don't overlap with what they want.} \\
\\ 
\texttt{Wait, the secret proposal is when we both submit our desired items. If our proposals are complementary, we get our items. So maybe I should propose taking items that the other player might not want. But how to know?} \\
\\ 
\texttt{Alternatively, since I value all items equally, maybe I should aim for a split where I get as many items as possible without conflicting. For example, if I take the bucket and battery, that's} \\
            }
        }
    }
    & \\ \\

    \theutterance \stepcounter{utterance}  
    & & & \multicolumn{4}{p{0.6\linewidth}}{
        \cellcolor[rgb]{0.9,0.9,0.9}{
            \makecell[{{p{\linewidth}}}]{
                \texttt{\tiny{[GM$\langle$P2]}}
                \texttt{<think>} \\
\texttt{Okay, let's see. I need to figure out how to divide the items between me and the other player. The items are 1 bucket, 3 knives, and 1 battery. My value function is 2 for each item, so each item is worth 2 points to me. The other player's values are secret, so I don't know what they prefer. The key is to make a proposal that's complementary to theirs so we both get points.} \\
\\ 
\texttt{Since I value all items equally, maybe I should aim for a fair split. But since the other player's preferences are unknown, I need to consider possible scenarios. Let's think about possible splits. If I take the bucket and battery, that's two items, leaving three knives. The other player could take all three knives, but that might not be complementary if they also wanted the bucket or battery. Alternatively, if I take some knives, maybe the other player can take the rest. But with three knives, splitting them 2{-}1 could work. But how to communicate this?} \\
\\ 
\texttt{I can send up to 5 messages. Maybe start by proposing a split. For example, I could suggest taking the bucket and battery, leaving the knives for the other player. But if they also value the bucket or battery highly, they might not agree. Alternatively, I could propose taking two knives and one of the other items. Wait, but my value is 2 per item, so each item is equally valuable. So maybe I should maximize the number of items I get. Since there are 5 items total, if I take 3 and they take 2, that's possible. But I need to make sure that their proposal doesn't overlap with mine.} \\
\\ 
\texttt{Wait, the secret proposal is when we both submit our desired items. If our proposals are complementary, we get our items. So maybe I should propose taking items that the other player might not want. But how to know? Since I don't know their values, I have to make an educated guess. The other player might value the bucket or battery more, or maybe the knives. If I take the bucket and battery, that's two items, and then take one knife, making three items total. Then the other player could take the remaining two knives. That would be complementary. Let me check: if I take 1 bucket, 1 battery, and 1 knife, that's 3 items. The other player would take 2} \\
            }
        }
    }
     \\ \\

    \theutterance \stepcounter{utterance}  
    & & & \multicolumn{2}{p{0.3\linewidth}}{
        \cellcolor[rgb]{0.9,0.9,0.9}{
            \makecell[{{p{\linewidth}}}]{
                \texttt{\tiny{[GM$|$GM]}}
                \texttt{continue} \\
            }
        }
    }
    & & \\ \\

    \theutterance \stepcounter{utterance}  
    & & & \multicolumn{2}{p{0.3\linewidth}}{
        \cellcolor[rgb]{0.9,0.9,0.9}{
            \makecell[{{p{\linewidth}}}]{
                \texttt{\tiny{[GM$|$GM]}}
                \texttt{<think>} \\
\texttt{Okay, let's see. I need to figure out how to divide the items between me and the other player. The items are 1 bucket, 3 knives, and 1 battery. My value function is 2 for each item, so each item is worth 2 points to me. The other player's values are secret, so I don't know what they prefer. The key is to make a proposal that's complementary to theirs so we both get points.} \\
\\ 
\texttt{Since I value all items equally, maybe I should aim for a fair split. But since the other player's preferences are unknown, I need to consider possible scenarios. Let's think about possible splits. If I take the bucket and battery, that's two items, leaving three knives. The other player could take all three knives, but that might not be complementary if they also wanted the bucket or battery. Alternatively, if I take some knives, maybe the other player can take the rest. But with three knives, splitting them 2{-}1 could work. But how to communicate this?} \\
\\ 
\texttt{I can send up to 5 messages. Maybe start by proposing a split. For example, I could suggest taking the bucket and battery, leaving the knives for the other player. But if they also value the bucket or battery highly, they might not agree. Alternatively, I could propose taking two knives and one of the other items. Wait, but my value is 2 per item, so each item is equally valuable. So maybe I should maximize the number of items I get. Since there are 5 items total, if I take 3 and they take 2, that's possible. But I need to make sure that their proposal doesn't overlap with mine.} \\
\\ 
\texttt{Wait, the secret proposal is when we both submit our desired items. If our proposals are complementary, we get our items. So maybe I should propose taking items that the other player might not want. But how to know? Since I don't know their values, I have to make an educated guess. The other player might value the bucket or battery more, or maybe the knives. If I take the bucket and battery, that's two items, and then take one knife, making three items total. Then the other player could take the remaining two knives. That would be complementary. Let me check: if I take 1 bucket, 1 battery, and 1 knife, that's 3 items. The other player would take 2} \\
            }
        }
    }
    & & \\ \\

    \theutterance \stepcounter{utterance}  
    & & \multicolumn{4}{p{0.6\linewidth}}{
        \cellcolor[rgb]{0.9,0.9,0.9}{
            \makecell[{{p{\linewidth}}}]{
                \texttt{\tiny{[P1$\langle$GM]}}
                \texttt{<think>} \\
\texttt{Okay, let's see. I need to figure out how to divide the items between me and the other player. The items are 1 bucket, 3 knives, and 1 battery. My value function is 2 for each item, so each item is worth 2 points to me. The other player's values are secret, so I don't know what they prefer. The key is to make a proposal that's complementary to theirs so we both get points.} \\
\\ 
\texttt{Since I value all items equally, maybe I should aim for a fair split. But since the other player's preferences are unknown, I need to consider possible scenarios. Let's think about possible splits. If I take the bucket and battery, that's two items, leaving three knives. The other player could take all three knives, but that might not be complementary if they also wanted the bucket or battery. Alternatively, if I take some knives, maybe the other player can take the rest. But with three knives, splitting them 2{-}1 could work. But how to communicate this?} \\
\\ 
\texttt{I can send up to 5 messages. Maybe start by proposing a split. For example, I could suggest taking the bucket and battery, leaving the knives for the other player. But if they also value the bucket or battery highly, they might not agree. Alternatively, I could propose taking two knives and one of the other items. Wait, but my value is 2 per item, so each item is equally valuable. So maybe I should maximize the number of items I get. Since there are 5 items total, if I take 3 and they take 2, that's possible. But I need to make sure that their proposal doesn't overlap with mine.} \\
\\ 
\texttt{Wait, the secret proposal is when we both submit our desired items. If our proposals are complementary, we get our items. So maybe I should propose taking items that the other player might not want. But how to know? Since I don't know their values, I have to make an educated guess. The other player might value the bucket or battery more, or maybe the knives. If I take the bucket and battery, that's two items, and then take one knife, making three items total. Then the other player could take the remaining two knives. That would be complementary. Let me check: if I take 1 bucket, 1 battery, and 1 knife, that's 3 items. The other player would take 2} \\
            }
        }
    }
    & \\ \\

    \theutterance \stepcounter{utterance}  
    & \multicolumn{4}{p{0.6\linewidth}}{
        \cellcolor[rgb]{0.9,0.9,0.9}{
            \makecell[{{p{\linewidth}}}]{
                \texttt{\tiny{[P1$\rangle$GM]}}
                \texttt{<think>} \\
\texttt{Okay, let's try to figure this out. The user is trying to divide the items between themselves and another player. The items are 1 bucket, 3 knives, and 1 battery. The user's value function is 2 for each item, so they value all items equally. The other player's values are secret, so the user doesn't know what they prefer.} \\
\\ 
\texttt{The user's goal is to maximize their own score, which is the sum of the values of the items they receive. If their proposal and the other player's proposal are complementary (i.e., they don't overlap and all items are distributed), then both get their scores. Otherwise, both get zero.} \\
\\ 
\texttt{In the previous attempt, the user considered taking the bucket, battery, and one knife, leaving two knives for the other player. That would give the user 3 items (3*2=6 points) and the other player 2 items (2*2=4 points). But the user is unsure if this is the best approach because they don't know the other player's preferences.} \\
\\ 
\texttt{Another possibility is to take more knives. Since there are three knives, maybe taking two knives and one other item. For example, two knives and the bucket. That would leave one knife and the battery for the other player. But again, the user doesn't know if the other player values the battery or bucket more.} \\
\\ 
\texttt{Wait, the user's value is the same for all items, so they just want as many items as possible. The maximum number of items they can get without overlapping with the other player's proposal is 3 (since there are 5 items total). So if the user takes 3 items, the other player can take 2. But the user needs to make sure that the other player's proposal doesn't conflict. Since the other player's values are unknown, the user has to make a proposal that is likely to be complementary regardless of the other player's preferences.} \\
\\ 
\texttt{If the user takes the bucket and battery (2 items), then they can take one knife, making 3 items. The other player would then have 2 knives left. Alternatively, if the user takes all three knives, that's 3 items, leaving the bucket and battery for the other player. But if the other player wants both the bucket and battery, that's two items, which would be complementary. However, if the other player also wants some knives, taking all three might not be complementary.} \\
            }
        }
    }
    & & \\ \\

    \theutterance \stepcounter{utterance}  
    & & & \multicolumn{2}{p{0.3\linewidth}}{
        \cellcolor[rgb]{0.9,0.9,0.9}{
            \makecell[{{p{\linewidth}}}]{
                \texttt{\tiny{[GM$|$GM]}}
                \texttt{continue} \\
            }
        }
    }
    & & \\ \\

    \theutterance \stepcounter{utterance}  
    & & & \multicolumn{2}{p{0.3\linewidth}}{
        \cellcolor[rgb]{0.9,0.9,0.9}{
            \makecell[{{p{\linewidth}}}]{
                \texttt{\tiny{[GM$|$GM]}}
                \texttt{<think>} \\
\texttt{Okay, let's try to figure this out. The user is trying to divide the items between themselves and another player. The items are 1 bucket, 3 knives, and 1 battery. The user's value function is 2 for each item, so they value all items equally. The other player's values are secret, so the user doesn't know what they prefer.} \\
\\ 
\texttt{The user's goal is to maximize their own score, which is the sum of the values of the items they receive. If their proposal and the other player's proposal are complementary (i.e., they don't overlap and all items are distributed), then both get their scores. Otherwise, both get zero.} \\
\\ 
\texttt{In the previous attempt, the user considered taking the bucket, battery, and one knife, leaving two knives for the other player. That would give the user 3 items (3*2=6 points) and the other player 2 items (2*2=4 points). But the user is unsure if this is the best approach because they don't know the other player's preferences.} \\
\\ 
\texttt{Another possibility is to take more knives. Since there are three knives, maybe taking two knives and one other item. For example, two knives and the bucket. That would leave one knife and the battery for the other player. But again, the user doesn't know if the other player values the battery or bucket more.} \\
\\ 
\texttt{Wait, the user's value is the same for all items, so they just want as many items as possible. The maximum number of items they can get without overlapping with the other player's proposal is 3 (since there are 5 items total). So if the user takes 3 items, the other player can take 2. But the user needs to make sure that the other player's proposal doesn't conflict. Since the other player's values are unknown, the user has to make a proposal that is likely to be complementary regardless of the other player's preferences.} \\
\\ 
\texttt{If the user takes the bucket and battery (2 items), then they can take one knife, making 3 items. The other player would then have 2 knives left. Alternatively, if the user takes all three knives, that's 3 items, leaving the bucket and battery for the other player. But if the other player wants both the bucket and battery, that's two items, which would be complementary. However, if the other player also wants some knives, taking all three might not be complementary.} \\
            }
        }
    }
    & & \\ \\

    \theutterance \stepcounter{utterance}  
    & & \multicolumn{4}{p{0.6\linewidth}}{
        \cellcolor[rgb]{0.9,0.9,0.9}{
            \makecell[{{p{\linewidth}}}]{
                \texttt{\tiny{[GM$\rangle$P2]}}
                \texttt{<think>} \\
\texttt{Okay, let's try to figure this out. The user is trying to divide the items between themselves and another player. The items are 1 bucket, 3 knives, and 1 battery. The user's value function is 2 for each item, so they value all items equally. The other player's values are secret, so the user doesn't know what they prefer.} \\
\\ 
\texttt{The user's goal is to maximize their own score, which is the sum of the values of the items they receive. If their proposal and the other player's proposal are complementary (i.e., they don't overlap and all items are distributed), then both get their scores. Otherwise, both get zero.} \\
\\ 
\texttt{In the previous attempt, the user considered taking the bucket, battery, and one knife, leaving two knives for the other player. That would give the user 3 items (3*2=6 points) and the other player 2 items (2*2=4 points). But the user is unsure if this is the best approach because they don't know the other player's preferences.} \\
\\ 
\texttt{Another possibility is to take more knives. Since there are three knives, maybe taking two knives and one other item. For example, two knives and the bucket. That would leave one knife and the battery for the other player. But again, the user doesn't know if the other player values the battery or bucket more.} \\
\\ 
\texttt{Wait, the user's value is the same for all items, so they just want as many items as possible. The maximum number of items they can get without overlapping with the other player's proposal is 3 (since there are 5 items total). So if the user takes 3 items, the other player can take 2. But the user needs to make sure that the other player's proposal doesn't conflict. Since the other player's values are unknown, the user has to make a proposal that is likely to be complementary regardless of the other player's preferences.} \\
\\ 
\texttt{If the user takes the bucket and battery (2 items), then they can take one knife, making 3 items. The other player would then have 2 knives left. Alternatively, if the user takes all three knives, that's 3 items, leaving the bucket and battery for the other player. But if the other player wants both the bucket and battery, that's two items, which would be complementary. However, if the other player also wants some knives, taking all three might not be complementary.} \\
            }
        }
    }
    & \\ \\

    \theutterance \stepcounter{utterance}  
    & & & \multicolumn{4}{p{0.6\linewidth}}{
        \cellcolor[rgb]{0.9,0.9,0.9}{
            \makecell[{{p{\linewidth}}}]{
                \texttt{\tiny{[GM$\langle$P2]}}
                \texttt{<think>} \\
\texttt{Okay, let's try to figure this out. The user is trying to divide the items between themselves and another player. The items are 1 bucket, 3 knives, and 1 battery. The user's value function is 2 for each item, so they value all items equally. The other player's values are secret, so the user doesn't know what they prefer.} \\
\\ 
\texttt{The user's goal is to maximize their own score, which is the sum of the values of the items they receive. If their proposal and the other player's proposal are complementary (i.e., they don't overlap and all items are distributed), then both get their scores. Otherwise, both get zero.} \\
\\ 
\texttt{In the previous attempt, the user considered taking the bucket, battery, and one knife, leaving two knives for the other player. That would give the user 3 items (3*2=6 points) and the other player 2 items (2*2=4 points). But the user is unsure if this is the best approach because they don't know the other player's preferences.} \\
\\ 
\texttt{Another possibility is to take more knives. Since there are three knives, maybe taking two knives and one other item. For example, two knives and the bucket. That would leave one knife and the battery for the other player. But again, the user doesn't know if the other player values the battery or bucket more.} \\
\\ 
\texttt{Wait, the user's value is the same for all items, so they just want as many items as possible. The maximum number of items they can get without overlapping with the other player's proposal is 3 (since there are 5 items total). So if the user takes 3 items, the other player can take 2. But the user needs to make sure that the other player's proposal doesn't conflict. Since the other player's values are unknown, the user has to make a proposal that is likely to be complementary regardless of the other player's preferences.} \\
\\ 
\texttt{If the user takes the bucket and battery (2 items), then they can take one knife, making 3 items. The other player would then have 2 knives left. Alternatively, if the user takes all three knives, that's 3 items, leaving the bucket and battery for the other player. But if the other player wants both the bucket and battery, that's two items, which would be complementary. However, if the other player also wants some knives, taking all three might not be complementary.} \\
            }
        }
    }
     \\ \\

    \theutterance \stepcounter{utterance}  
    & & & \multicolumn{2}{p{0.3\linewidth}}{
        \cellcolor[rgb]{0.9,0.9,0.9}{
            \makecell[{{p{\linewidth}}}]{
                \texttt{\tiny{[GM$|$GM]}}
                \texttt{continue} \\
            }
        }
    }
    & & \\ \\

    \theutterance \stepcounter{utterance}  
    & & & \multicolumn{2}{p{0.3\linewidth}}{
        \cellcolor[rgb]{0.9,0.9,0.9}{
            \makecell[{{p{\linewidth}}}]{
                \texttt{\tiny{[GM$|$GM]}}
                \texttt{<think>} \\
\texttt{Okay, let's try to figure this out. The user is trying to divide the items between themselves and another player. The items are 1 bucket, 3 knives, and 1 battery. The user's value function is 2 for each item, so they value all items equally. The other player's values are secret, so the user doesn't know what they prefer.} \\
\\ 
\texttt{The user's goal is to maximize their own score, which is the sum of the values of the items they receive. If their proposal and the other player's proposal are complementary (i.e., they don't overlap and all items are distributed), then both get their scores. Otherwise, both get zero.} \\
\\ 
\texttt{In the previous attempt, the user considered taking the bucket, battery, and one knife, leaving two knives for the other player. That would give the user 3 items (3*2=6 points) and the other player 2 items (2*2=4 points). But the user is unsure if this is the best approach because they don't know the other player's preferences.} \\
\\ 
\texttt{Another possibility is to take more knives. Since there are three knives, maybe taking two knives and one other item. For example, two knives and the bucket. That would leave one knife and the battery for the other player. But again, the user doesn't know if the other player values the battery or bucket more.} \\
\\ 
\texttt{Wait, the user's value is the same for all items, so they just want as many items as possible. The maximum number of items they can get without overlapping with the other player's proposal is 3 (since there are 5 items total). So if the user takes 3 items, the other player can take 2. But the user needs to make sure that the other player's proposal doesn't conflict. Since the other player's values are unknown, the user has to make a proposal that is likely to be complementary regardless of the other player's preferences.} \\
\\ 
\texttt{If the user takes the bucket and battery (2 items), then they can take one knife, making 3 items. The other player would then have 2 knives left. Alternatively, if the user takes all three knives, that's 3 items, leaving the bucket and battery for the other player. But if the other player wants both the bucket and battery, that's two items, which would be complementary. However, if the other player also wants some knives, taking all three might not be complementary.} \\
            }
        }
    }
    & & \\ \\

    \theutterance \stepcounter{utterance}  
    & & \multicolumn{4}{p{0.6\linewidth}}{
        \cellcolor[rgb]{0.9,0.9,0.9}{
            \makecell[{{p{\linewidth}}}]{
                \texttt{\tiny{[P1$\langle$GM]}}
                \texttt{<think>} \\
\texttt{Okay, let's try to figure this out. The user is trying to divide the items between themselves and another player. The items are 1 bucket, 3 knives, and 1 battery. The user's value function is 2 for each item, so they value all items equally. The other player's values are secret, so the user doesn't know what they prefer.} \\
\\ 
\texttt{The user's goal is to maximize their own score, which is the sum of the values of the items they receive. If their proposal and the other player's proposal are complementary (i.e., they don't overlap and all items are distributed), then both get their scores. Otherwise, both get zero.} \\
\\ 
\texttt{In the previous attempt, the user considered taking the bucket, battery, and one knife, leaving two knives for the other player. That would give the user 3 items (3*2=6 points) and the other player 2 items (2*2=4 points). But the user is unsure if this is the best approach because they don't know the other player's preferences.} \\
\\ 
\texttt{Another possibility is to take more knives. Since there are three knives, maybe taking two knives and one other item. For example, two knives and the bucket. That would leave one knife and the battery for the other player. But again, the user doesn't know if the other player values the battery or bucket more.} \\
\\ 
\texttt{Wait, the user's value is the same for all items, so they just want as many items as possible. The maximum number of items they can get without overlapping with the other player's proposal is 3 (since there are 5 items total). So if the user takes 3 items, the other player can take 2. But the user needs to make sure that the other player's proposal doesn't conflict. Since the other player's values are unknown, the user has to make a proposal that is likely to be complementary regardless of the other player's preferences.} \\
\\ 
\texttt{If the user takes the bucket and battery (2 items), then they can take one knife, making 3 items. The other player would then have 2 knives left. Alternatively, if the user takes all three knives, that's 3 items, leaving the bucket and battery for the other player. But if the other player wants both the bucket and battery, that's two items, which would be complementary. However, if the other player also wants some knives, taking all three might not be complementary.} \\
            }
        }
    }
    & \\ \\

    \theutterance \stepcounter{utterance}  
    & \multicolumn{4}{p{0.6\linewidth}}{
        \cellcolor[rgb]{0.9,0.9,0.9}{
            \makecell[{{p{\linewidth}}}]{
                \texttt{\tiny{[P1$\rangle$GM]}}
                \texttt{<think>} \\
\texttt{Okay, let's try to figure this out. The user is trying to divide the items between themselves and another player. The items are 1 bucket, 3 knives, and 1 battery. The user's value function is 2 for each item, so they value all items equally. The other player's values are secret, so the user doesn't know what they prefer.} \\
\\ 
\texttt{The user's goal is to maximize their own score, which is the sum of the values of the items they receive. If their proposal and the other player's proposal are complementary (i.e., they don't overlap and all items are distributed), then both get their scores. Otherwise, both get zero.} \\
\\ 
\texttt{In the previous attempt, the user considered taking the bucket, battery, and one knife, leaving two knives for the other player. That would give the user 3 items (3*2=6 points) and the other player 2 items (2*2=4 points). But the user is unsure if this is the best approach because they don't know the other player's preferences.} \\
\\ 
\texttt{Another possibility is to take more knives. Since there are three knives, maybe taking two knives and one other item. For example, two knives and the bucket. That would leave one knife and the battery for the other player. But again, the user doesn't know if the other player values the battery or bucket more.} \\
\\ 
\texttt{Wait, the user's value is the same for all items, so they just want as many items as possible. The maximum number of items they can get without overlapping with the other player's proposal is 3 (since there are 5 items total). So if the user takes 3 items, the other player can take 2. But the user needs to make sure that the other player's proposal doesn't conflict. Since the other player's values are unknown, the user has to make a proposal that is likely to be complementary regardless of the other player's preferences.} \\
\\ 
\texttt{If the user takes the bucket and battery (2 items), then they can take one knife, making 3 items. The other player would then have 2 knives left. Alternatively, if the user takes all three knives, that's 3 items, leaving the bucket and battery for the other player. But if the other player wants both the bucket and battery, that's two items, which would be complementary. However, if the other player also wants some knives, taking all three might not be complementary.} \\
            }
        }
    }
    & & \\ \\

    \theutterance \stepcounter{utterance}  
    & & & \multicolumn{2}{p{0.3\linewidth}}{
        \cellcolor[rgb]{0.9,0.9,0.9}{
            \makecell[{{p{\linewidth}}}]{
                \texttt{\tiny{[GM$|$GM]}}
                \texttt{continue} \\
            }
        }
    }
    & & \\ \\

    \theutterance \stepcounter{utterance}  
    & & & \multicolumn{2}{p{0.3\linewidth}}{
        \cellcolor[rgb]{0.9,0.9,0.9}{
            \makecell[{{p{\linewidth}}}]{
                \texttt{\tiny{[GM$|$GM]}}
                \texttt{<think>} \\
\texttt{Okay, let's try to figure this out. The user is trying to divide the items between themselves and another player. The items are 1 bucket, 3 knives, and 1 battery. The user's value function is 2 for each item, so they value all items equally. The other player's values are secret, so the user doesn't know what they prefer.} \\
\\ 
\texttt{The user's goal is to maximize their own score, which is the sum of the values of the items they receive. If their proposal and the other player's proposal are complementary (i.e., they don't overlap and all items are distributed), then both get their scores. Otherwise, both get zero.} \\
\\ 
\texttt{In the previous attempt, the user considered taking the bucket, battery, and one knife, leaving two knives for the other player. That would give the user 3 items (3*2=6 points) and the other player 2 items (2*2=4 points). But the user is unsure if this is the best approach because they don't know the other player's preferences.} \\
\\ 
\texttt{Another possibility is to take more knives. Since there are three knives, maybe taking two knives and one other item. For example, two knives and the bucket. That would leave one knife and the battery for the other player. But again, the user doesn't know if the other player values the battery or bucket more.} \\
\\ 
\texttt{Wait, the user's value is the same for all items, so they just want as many items as possible. The maximum number of items they can get without overlapping with the other player's proposal is 3 (since there are 5 items total). So if the user takes 3 items, the other player can take 2. But the user needs to make sure that the other player's proposal doesn't conflict. Since the other player's values are unknown, the user has to make a proposal that is likely to be complementary regardless of the other player's preferences.} \\
\\ 
\texttt{If the user takes the bucket and battery (2 items), then they can take one knife, making 3 items. The other player would then have 2 knives left. Alternatively, if the user takes all three knives, that's 3 items, leaving the bucket and battery for the other player. But if the other player wants both the bucket and battery, that's two items, which would be complementary. However, if the other player also wants some knives, taking all three might not be complementary.} \\
            }
        }
    }
    & & \\ \\

    \theutterance \stepcounter{utterance}  
    & & \multicolumn{4}{p{0.6\linewidth}}{
        \cellcolor[rgb]{0.9,0.9,0.9}{
            \makecell[{{p{\linewidth}}}]{
                \texttt{\tiny{[GM$\rangle$P2]}}
                \texttt{<think>} \\
\texttt{Okay, let's try to figure this out. The user is trying to divide the items between themselves and another player. The items are 1 bucket, 3 knives, and 1 battery. The user's value function is 2 for each item, so they value all items equally. The other player's values are secret, so the user doesn't know what they prefer.} \\
\\ 
\texttt{The user's goal is to maximize their own score, which is the sum of the values of the items they receive. If their proposal and the other player's proposal are complementary (i.e., they don't overlap and all items are distributed), then both get their scores. Otherwise, both get zero.} \\
\\ 
\texttt{In the previous attempt, the user considered taking the bucket, battery, and one knife, leaving two knives for the other player. That would give the user 3 items (3*2=6 points) and the other player 2 items (2*2=4 points). But the user is unsure if this is the best approach because they don't know the other player's preferences.} \\
\\ 
\texttt{Another possibility is to take more knives. Since there are three knives, maybe taking two knives and one other item. For example, two knives and the bucket. That would leave one knife and the battery for the other player. But again, the user doesn't know if the other player values the battery or bucket more.} \\
\\ 
\texttt{Wait, the user's value is the same for all items, so they just want as many items as possible. The maximum number of items they can get without overlapping with the other player's proposal is 3 (since there are 5 items total). So if the user takes 3 items, the other player can take 2. But the user needs to make sure that the other player's proposal doesn't conflict. Since the other player's values are unknown, the user has to make a proposal that is likely to be complementary regardless of the other player's preferences.} \\
\\ 
\texttt{If the user takes the bucket and battery (2 items), then they can take one knife, making 3 items. The other player would then have 2 knives left. Alternatively, if the user takes all three knives, that's 3 items, leaving the bucket and battery for the other player. But if the other player wants both the bucket and battery, that's two items, which would be complementary. However, if the other player also wants some knives, taking all three might not be complementary.} \\
            }
        }
    }
    & \\ \\

    \theutterance \stepcounter{utterance}  
    & & & \multicolumn{4}{p{0.6\linewidth}}{
        \cellcolor[rgb]{0.9,0.9,0.9}{
            \makecell[{{p{\linewidth}}}]{
                \texttt{\tiny{[GM$\langle$P2]}}
                \texttt{<think>} \\
\texttt{Okay, let's try to figure this out. The user is trying to divide the items between themselves and another player. The items are 1 bucket, 3 knives, and 1 battery. The user's value function is 2 for each item, so they value all items equally. The other player's values are secret, so the user doesn't know what they prefer.} \\
\\ 
\texttt{The user's goal is to maximize their own score, which is the sum of the values of the items they receive. If their proposal and the other player's proposal are complementary (i.e., they don't overlap and all items are distributed), then both get their scores. Otherwise, both get zero.} \\
\\ 
\texttt{In the previous attempt, the user considered taking the bucket, battery, and one knife, leaving two knives for the other player. That would give the user 3 items (3*2=6 points) and the other player 2 items (2*2=4 points). But the user is unsure if this is the best approach because they don't know the other player's preferences.} \\
\\ 
\texttt{Another possibility is to take more knives. Since there are three knives, maybe taking two knives and one other item. For example, two knives and the bucket. That would leave one knife and the battery for the other player. But again, the user doesn't know if the other player values the battery or bucket more.} \\
\\ 
\texttt{Wait, the user's value is the same for all items, so they just want as many items as possible. The maximum number of items they can get without overlapping with the other player's proposal is 3 (since there are 5 items total). So if the user takes 3 items, the other player can take 2. But the user needs to make sure that the other player's proposal doesn't conflict. Since the other player's values are unknown, the user has to make a proposal that is likely to be complementary regardless of the other player's preferences.} \\
\\ 
\texttt{If the user takes the bucket and battery (2 items), then they can take one knife, making 3 items. The other player would then have 2 knives left. Alternatively, if the user takes all three knives, that's 3 items, leaving the bucket and battery for the other player. But if the other player wants both the bucket and battery, that's two items, which would be complementary. However, if the other player also wants some knives, taking all three might not be complementary.} \\
            }
        }
    }
     \\ \\

    \theutterance \stepcounter{utterance}  
    & & & \multicolumn{2}{p{0.3\linewidth}}{
        \cellcolor[rgb]{0.9,0.9,0.9}{
            \makecell[{{p{\linewidth}}}]{
                \texttt{\tiny{[GM$|$GM]}}
                \texttt{continue} \\
            }
        }
    }
    & & \\ \\

    \theutterance \stepcounter{utterance}  
    & & & \multicolumn{2}{p{0.3\linewidth}}{
        \cellcolor[rgb]{0.9,0.9,0.9}{
            \makecell[{{p{\linewidth}}}]{
                \texttt{\tiny{[GM$|$GM]}}
                \texttt{<think>} \\
\texttt{Okay, let's try to figure this out. The user is trying to divide the items between themselves and another player. The items are 1 bucket, 3 knives, and 1 battery. The user's value function is 2 for each item, so they value all items equally. The other player's values are secret, so the user doesn't know what they prefer.} \\
\\ 
\texttt{The user's goal is to maximize their own score, which is the sum of the values of the items they receive. If their proposal and the other player's proposal are complementary (i.e., they don't overlap and all items are distributed), then both get their scores. Otherwise, both get zero.} \\
\\ 
\texttt{In the previous attempt, the user considered taking the bucket, battery, and one knife, leaving two knives for the other player. That would give the user 3 items (3*2=6 points) and the other player 2 items (2*2=4 points). But the user is unsure if this is the best approach because they don't know the other player's preferences.} \\
\\ 
\texttt{Another possibility is to take more knives. Since there are three knives, maybe taking two knives and one other item. For example, two knives and the bucket. That would leave one knife and the battery for the other player. But again, the user doesn't know if the other player values the battery or bucket more.} \\
\\ 
\texttt{Wait, the user's value is the same for all items, so they just want as many items as possible. The maximum number of items they can get without overlapping with the other player's proposal is 3 (since there are 5 items total). So if the user takes 3 items, the other player can take 2. But the user needs to make sure that the other player's proposal doesn't conflict. Since the other player's values are unknown, the user has to make a proposal that is likely to be complementary regardless of the other player's preferences.} \\
\\ 
\texttt{If the user takes the bucket and battery (2 items), then they can take one knife, making 3 items. The other player would then have 2 knives left. Alternatively, if the user takes all three knives, that's 3 items, leaving the bucket and battery for the other player. But if the other player wants both the bucket and battery, that's two items, which would be complementary. However, if the other player also wants some knives, taking all three might not be complementary.} \\
            }
        }
    }
    & & \\ \\

    \theutterance \stepcounter{utterance}  
    & & \multicolumn{4}{p{0.6\linewidth}}{
        \cellcolor[rgb]{0.9,0.9,0.9}{
            \makecell[{{p{\linewidth}}}]{
                \texttt{\tiny{[P1$\langle$GM]}}
                \texttt{<think>} \\
\texttt{Okay, let's try to figure this out. The user is trying to divide the items between themselves and another player. The items are 1 bucket, 3 knives, and 1 battery. The user's value function is 2 for each item, so they value all items equally. The other player's values are secret, so the user doesn't know what they prefer.} \\
\\ 
\texttt{The user's goal is to maximize their own score, which is the sum of the values of the items they receive. If their proposal and the other player's proposal are complementary (i.e., they don't overlap and all items are distributed), then both get their scores. Otherwise, both get zero.} \\
\\ 
\texttt{In the previous attempt, the user considered taking the bucket, battery, and one knife, leaving two knives for the other player. That would give the user 3 items (3*2=6 points) and the other player 2 items (2*2=4 points). But the user is unsure if this is the best approach because they don't know the other player's preferences.} \\
\\ 
\texttt{Another possibility is to take more knives. Since there are three knives, maybe taking two knives and one other item. For example, two knives and the bucket. That would leave one knife and the battery for the other player. But again, the user doesn't know if the other player values the battery or bucket more.} \\
\\ 
\texttt{Wait, the user's value is the same for all items, so they just want as many items as possible. The maximum number of items they can get without overlapping with the other player's proposal is 3 (since there are 5 items total). So if the user takes 3 items, the other player can take 2. But the user needs to make sure that the other player's proposal doesn't conflict. Since the other player's values are unknown, the user has to make a proposal that is likely to be complementary regardless of the other player's preferences.} \\
\\ 
\texttt{If the user takes the bucket and battery (2 items), then they can take one knife, making 3 items. The other player would then have 2 knives left. Alternatively, if the user takes all three knives, that's 3 items, leaving the bucket and battery for the other player. But if the other player wants both the bucket and battery, that's two items, which would be complementary. However, if the other player also wants some knives, taking all three might not be complementary.} \\
            }
        }
    }
    & \\ \\

    \theutterance \stepcounter{utterance}  
    & \multicolumn{4}{p{0.6\linewidth}}{
        \cellcolor[rgb]{0.9,0.9,0.9}{
            \makecell[{{p{\linewidth}}}]{
                \texttt{\tiny{[P1$\rangle$GM]}}
                \texttt{<think>} \\
\texttt{Okay, let's try to figure this out. The user is trying to divide the items between themselves and another player. The items are 1 bucket, 3 knives, and 1 battery. The user's value function is 2 for each item, so they value all items equally. The other player's values are secret, so the user doesn't know what they prefer.} \\
\\ 
\texttt{The user's goal is to maximize their own score, which is the sum of the values of the items they receive. If their proposal and the other player's proposal are complementary (i.e., they don't overlap and all items are distributed), then both get their scores. Otherwise, both get zero.} \\
\\ 
\texttt{In the previous attempt, the user considered taking the bucket, battery, and one knife, leaving two knives for the other player. That would give the user 3 items (3*2=6 points) and the other player 2 items (2*2=4 points). But the user is unsure if this is the best approach because they don't know the other player's preferences.} \\
\\ 
\texttt{Another possibility is to take more knives. Since there are three knives, maybe taking two knives and one other item. For example, two knives and the bucket. That would leave one knife and the battery for the other player. But again, the user doesn't know if the other player values the battery or bucket more.} \\
\\ 
\texttt{Wait, the user's value is the same for all items, so they just want as many items as possible. The maximum number of items they can get without overlapping with the other player's proposal is 3 (since there are 5 items total). So if the user takes 3 items, the other player can take 2. But the user needs to make sure that the other player's proposal doesn't conflict. Since the other player's values are unknown, the user has to make a proposal that is likely to be complementary regardless of the other player's preferences.} \\
\\ 
\texttt{If the user takes the bucket and battery (2 items), then they can take one knife, making 3 items. The other player would then have 2 knives left. Alternatively, if the user takes all three knives, that's 3 items, leaving the bucket and battery for the other player. But if the other player wants both the bucket and battery, that's two items, which would be complementary. However, if the other player also wants some knives, taking all three might not be complementary.} \\
            }
        }
    }
    & & \\ \\

    \theutterance \stepcounter{utterance}  
    & & & \multicolumn{2}{p{0.3\linewidth}}{
        \cellcolor[rgb]{0.9,0.9,0.9}{
            \makecell[{{p{\linewidth}}}]{
                \texttt{\tiny{[GM$|$GM]}}
                \texttt{continue} \\
            }
        }
    }
    & & \\ \\

    \theutterance \stepcounter{utterance}  
    & & & \multicolumn{2}{p{0.3\linewidth}}{
        \cellcolor[rgb]{0.9,0.9,0.9}{
            \makecell[{{p{\linewidth}}}]{
                \texttt{\tiny{[GM$|$GM]}}
                \texttt{<think>} \\
\texttt{Okay, let's try to figure this out. The user is trying to divide the items between themselves and another player. The items are 1 bucket, 3 knives, and 1 battery. The user's value function is 2 for each item, so they value all items equally. The other player's values are secret, so the user doesn't know what they prefer.} \\
\\ 
\texttt{The user's goal is to maximize their own score, which is the sum of the values of the items they receive. If their proposal and the other player's proposal are complementary (i.e., they don't overlap and all items are distributed), then both get their scores. Otherwise, both get zero.} \\
\\ 
\texttt{In the previous attempt, the user considered taking the bucket, battery, and one knife, leaving two knives for the other player. That would give the user 3 items (3*2=6 points) and the other player 2 items (2*2=4 points). But the user is unsure if this is the best approach because they don't know the other player's preferences.} \\
\\ 
\texttt{Another possibility is to take more knives. Since there are three knives, maybe taking two knives and one other item. For example, two knives and the bucket. That would leave one knife and the battery for the other player. But again, the user doesn't know if the other player values the battery or bucket more.} \\
\\ 
\texttt{Wait, the user's value is the same for all items, so they just want as many items as possible. The maximum number of items they can get without overlapping with the other player's proposal is 3 (since there are 5 items total). So if the user takes 3 items, the other player can take 2. But the user needs to make sure that the other player's proposal doesn't conflict. Since the other player's values are unknown, the user has to make a proposal that is likely to be complementary regardless of the other player's preferences.} \\
\\ 
\texttt{If the user takes the bucket and battery (2 items), then they can take one knife, making 3 items. The other player would then have 2 knives left. Alternatively, if the user takes all three knives, that's 3 items, leaving the bucket and battery for the other player. But if the other player wants both the bucket and battery, that's two items, which would be complementary. However, if the other player also wants some knives, taking all three might not be complementary.} \\
            }
        }
    }
    & & \\ \\

    \theutterance \stepcounter{utterance}  
    & & \multicolumn{4}{p{0.6\linewidth}}{
        \cellcolor[rgb]{0.9,0.9,0.9}{
            \makecell[{{p{\linewidth}}}]{
                \texttt{\tiny{[GM$\rangle$P2]}}
                \texttt{<think>} \\
\texttt{Okay, let's try to figure this out. The user is trying to divide the items between themselves and another player. The items are 1 bucket, 3 knives, and 1 battery. The user's value function is 2 for each item, so they value all items equally. The other player's values are secret, so the user doesn't know what they prefer.} \\
\\ 
\texttt{The user's goal is to maximize their own score, which is the sum of the values of the items they receive. If their proposal and the other player's proposal are complementary (i.e., they don't overlap and all items are distributed), then both get their scores. Otherwise, both get zero.} \\
\\ 
\texttt{In the previous attempt, the user considered taking the bucket, battery, and one knife, leaving two knives for the other player. That would give the user 3 items (3*2=6 points) and the other player 2 items (2*2=4 points). But the user is unsure if this is the best approach because they don't know the other player's preferences.} \\
\\ 
\texttt{Another possibility is to take more knives. Since there are three knives, maybe taking two knives and one other item. For example, two knives and the bucket. That would leave one knife and the battery for the other player. But again, the user doesn't know if the other player values the battery or bucket more.} \\
\\ 
\texttt{Wait, the user's value is the same for all items, so they just want as many items as possible. The maximum number of items they can get without overlapping with the other player's proposal is 3 (since there are 5 items total). So if the user takes 3 items, the other player can take 2. But the user needs to make sure that the other player's proposal doesn't conflict. Since the other player's values are unknown, the user has to make a proposal that is likely to be complementary regardless of the other player's preferences.} \\
\\ 
\texttt{If the user takes the bucket and battery (2 items), then they can take one knife, making 3 items. The other player would then have 2 knives left. Alternatively, if the user takes all three knives, that's 3 items, leaving the bucket and battery for the other player. But if the other player wants both the bucket and battery, that's two items, which would be complementary. However, if the other player also wants some knives, taking all three might not be complementary.} \\
            }
        }
    }
    & \\ \\

    \theutterance \stepcounter{utterance}  
    & & & \multicolumn{4}{p{0.6\linewidth}}{
        \cellcolor[rgb]{0.9,0.9,0.9}{
            \makecell[{{p{\linewidth}}}]{
                \texttt{\tiny{[GM$\langle$P2]}}
                \texttt{<think>} \\
\texttt{Okay, let's try to figure this out. The user is trying to divide the items between themselves and another player. The items are 1 bucket, 3 knives, and 1 battery. The user's value function is 2 for each item, so they value all items equally. The other player's values are secret, so the user doesn't know what they prefer.} \\
\\ 
\texttt{The user's goal is to maximize their own score, which is the sum of the values of the items they receive. If their proposal and the other player's proposal are complementary (i.e., they don't overlap and all items are distributed), then both get their scores. Otherwise, both get zero.} \\
\\ 
\texttt{In the previous attempt, the user considered taking the bucket, battery, and one knife, leaving two knives for the other player. That would give the user 3 items (3*2=6 points) and the other player 2 items (2*2=4 points). But the user is unsure if this is the best approach because they don't know the other player's preferences.} \\
\\ 
\texttt{Another possibility is to take more knives. Since there are three knives, maybe taking two knives and one other item. For example, two knives and the bucket. That would leave one knife and the battery for the other player. But again, the user doesn't know if the other player values the battery or bucket more.} \\
\\ 
\texttt{Wait, the user's value is the same for all items, so they just want as many items as possible. The maximum number of items they can get without overlapping with the other player's proposal is 3 (since there are 5 items total). So if the user takes 3 items, the other player can take 2. But the user needs to make sure that the other player's proposal doesn't conflict. Since the other player's values are unknown, the user has to make a proposal that is likely to be complementary regardless of the other player's preferences.} \\
\\ 
\texttt{If the user takes the bucket and battery (2 items), then they can take one knife, making 3 items. The other player would then have 2 knives left. Alternatively, if the user takes all three knives, that's 3 items, leaving the bucket and battery for the other player. But if the other player wants both the bucket and battery, that's two items, which would be complementary. However, if the other player also wants some knives, taking all three might not be complementary.} \\
            }
        }
    }
     \\ \\

    \theutterance \stepcounter{utterance}  
    & & & \multicolumn{2}{p{0.3\linewidth}}{
        \cellcolor[rgb]{0.9,0.9,0.9}{
            \makecell[{{p{\linewidth}}}]{
                \texttt{\tiny{[GM$|$GM]}}
                \texttt{continue} \\
            }
        }
    }
    & & \\ \\

    \theutterance \stepcounter{utterance}  
    & & & \multicolumn{2}{p{0.3\linewidth}}{
        \cellcolor[rgb]{0.9,0.9,0.9}{
            \makecell[{{p{\linewidth}}}]{
                \texttt{\tiny{[GM$|$GM]}}
                \texttt{<think>} \\
\texttt{Okay, let's try to figure this out. The user is trying to divide the items between themselves and another player. The items are 1 bucket, 3 knives, and 1 battery. The user's value function is 2 for each item, so they value all items equally. The other player's values are secret, so the user doesn't know what they prefer.} \\
\\ 
\texttt{The user's goal is to maximize their own score, which is the sum of the values of the items they receive. If their proposal and the other player's proposal are complementary (i.e., they don't overlap and all items are distributed), then both get their scores. Otherwise, both get zero.} \\
\\ 
\texttt{In the previous attempt, the user considered taking the bucket, battery, and one knife, leaving two knives for the other player. That would give the user 3 items (3*2=6 points) and the other player 2 items (2*2=4 points). But the user is unsure if this is the best approach because they don't know the other player's preferences.} \\
\\ 
\texttt{Another possibility is to take more knives. Since there are three knives, maybe taking two knives and one other item. For example, two knives and the bucket. That would leave one knife and the battery for the other player. But again, the user doesn't know if the other player values the battery or bucket more.} \\
\\ 
\texttt{Wait, the user's value is the same for all items, so they just want as many items as possible. The maximum number of items they can get without overlapping with the other player's proposal is 3 (since there are 5 items total). So if the user takes 3 items, the other player can take 2. But the user needs to make sure that the other player's proposal doesn't conflict. Since the other player's values are unknown, the user has to make a proposal that is likely to be complementary regardless of the other player's preferences.} \\
\\ 
\texttt{If the user takes the bucket and battery (2 items), then they can take one knife, making 3 items. The other player would then have 2 knives left. Alternatively, if the user takes all three knives, that's 3 items, leaving the bucket and battery for the other player. But if the other player wants both the bucket and battery, that's two items, which would be complementary. However, if the other player also wants some knives, taking all three might not be complementary.} \\
            }
        }
    }
    & & \\ \\

    \theutterance \stepcounter{utterance}  
    & & \multicolumn{4}{p{0.6\linewidth}}{
        \cellcolor[rgb]{0.9,0.9,0.9}{
            \makecell[{{p{\linewidth}}}]{
                \texttt{\tiny{[P1$\langle$GM]}}
                \texttt{<think>} \\
\texttt{Okay, let's try to figure this out. The user is trying to divide the items between themselves and another player. The items are 1 bucket, 3 knives, and 1 battery. The user's value function is 2 for each item, so they value all items equally. The other player's values are secret, so the user doesn't know what they prefer.} \\
\\ 
\texttt{The user's goal is to maximize their own score, which is the sum of the values of the items they receive. If their proposal and the other player's proposal are complementary (i.e., they don't overlap and all items are distributed), then both get their scores. Otherwise, both get zero.} \\
\\ 
\texttt{In the previous attempt, the user considered taking the bucket, battery, and one knife, leaving two knives for the other player. That would give the user 3 items (3*2=6 points) and the other player 2 items (2*2=4 points). But the user is unsure if this is the best approach because they don't know the other player's preferences.} \\
\\ 
\texttt{Another possibility is to take more knives. Since there are three knives, maybe taking two knives and one other item. For example, two knives and the bucket. That would leave one knife and the battery for the other player. But again, the user doesn't know if the other player values the battery or bucket more.} \\
\\ 
\texttt{Wait, the user's value is the same for all items, so they just want as many items as possible. The maximum number of items they can get without overlapping with the other player's proposal is 3 (since there are 5 items total). So if the user takes 3 items, the other player can take 2. But the user needs to make sure that the other player's proposal doesn't conflict. Since the other player's values are unknown, the user has to make a proposal that is likely to be complementary regardless of the other player's preferences.} \\
\\ 
\texttt{If the user takes the bucket and battery (2 items), then they can take one knife, making 3 items. The other player would then have 2 knives left. Alternatively, if the user takes all three knives, that's 3 items, leaving the bucket and battery for the other player. But if the other player wants both the bucket and battery, that's two items, which would be complementary. However, if the other player also wants some knives, taking all three might not be complementary.} \\
            }
        }
    }
    & \\ \\

    \theutterance \stepcounter{utterance}  
    & \multicolumn{4}{p{0.6\linewidth}}{
        \cellcolor[rgb]{0.9,0.9,0.9}{
            \makecell[{{p{\linewidth}}}]{
                \texttt{\tiny{[P1$\rangle$GM]}}
                \texttt{<think>} \\
\texttt{Okay, let's try to figure this out. The user is trying to divide the items between themselves and another player. The items are 1 bucket, 3 knives, and 1 battery. The user's value function is 2 for each item, so they value all items equally. The other player's values are secret, so the user doesn't know what they prefer.} \\
\\ 
\texttt{The user's goal is to maximize their own score, which is the sum of the values of the items they receive. If their proposal and the other player's proposal are complementary (i.e., they don't overlap and all items are distributed), then both get their scores. Otherwise, both get zero.} \\
\\ 
\texttt{In the previous attempt, the user considered taking the bucket, battery, and one knife, leaving two knives for the other player. That would give the user 3 items (3*2=6 points) and the other player 2 items (2*2=4 points). But the user is unsure if this is the best approach because they don't know the other player's preferences.} \\
\\ 
\texttt{Another possibility is to take more knives. Since there are three knives, maybe taking two knives and one other item. For example, two knives and the bucket. That would leave one knife and the battery for the other player. But again, the user doesn't know if the other player values the battery or bucket more.} \\
\\ 
\texttt{Wait, the user's value is the same for all items, so they just want as many items as possible. The maximum number of items they can get without overlapping with the other player's proposal is 3 (since there are 5 items total). So if the user takes 3 items, the other player can take 2. But the user needs to make sure that the other player's proposal doesn't conflict. Since the other player's values are unknown, the user has to make a proposal that is likely to be complementary regardless of the other player's preferences.} \\
\\ 
\texttt{If the user takes the bucket and battery (2 items), then they can take one knife, making 3 items. The other player would then have 2 knives left. Alternatively, if the user takes all three knives, that's 3 items, leaving the bucket and battery for the other player. But if the other player wants both the bucket and battery, that's two items, which would be complementary. However, if the other player also wants some knives, taking all three might not be complementary.} \\
            }
        }
    }
    & & \\ \\

    \theutterance \stepcounter{utterance}  
    & & & \multicolumn{2}{p{0.3\linewidth}}{
        \cellcolor[rgb]{0.9,0.9,0.9}{
            \makecell[{{p{\linewidth}}}]{
                \texttt{\tiny{[GM$|$GM]}}
                \texttt{continue} \\
            }
        }
    }
    & & \\ \\

    \theutterance \stepcounter{utterance}  
    & & & \multicolumn{2}{p{0.3\linewidth}}{
        \cellcolor[rgb]{0.9,0.9,0.9}{
            \makecell[{{p{\linewidth}}}]{
                \texttt{\tiny{[GM$|$GM]}}
                \texttt{<think>} \\
\texttt{Okay, let's try to figure this out. The user is trying to divide the items between themselves and another player. The items are 1 bucket, 3 knives, and 1 battery. The user's value function is 2 for each item, so they value all items equally. The other player's values are secret, so the user doesn't know what they prefer.} \\
\\ 
\texttt{The user's goal is to maximize their own score, which is the sum of the values of the items they receive. If their proposal and the other player's proposal are complementary (i.e., they don't overlap and all items are distributed), then both get their scores. Otherwise, both get zero.} \\
\\ 
\texttt{In the previous attempt, the user considered taking the bucket, battery, and one knife, leaving two knives for the other player. That would give the user 3 items (3*2=6 points) and the other player 2 items (2*2=4 points). But the user is unsure if this is the best approach because they don't know the other player's preferences.} \\
\\ 
\texttt{Another possibility is to take more knives. Since there are three knives, maybe taking two knives and one other item. For example, two knives and the bucket. That would leave one knife and the battery for the other player. But again, the user doesn't know if the other player values the battery or bucket more.} \\
\\ 
\texttt{Wait, the user's value is the same for all items, so they just want as many items as possible. The maximum number of items they can get without overlapping with the other player's proposal is 3 (since there are 5 items total). So if the user takes 3 items, the other player can take 2. But the user needs to make sure that the other player's proposal doesn't conflict. Since the other player's values are unknown, the user has to make a proposal that is likely to be complementary regardless of the other player's preferences.} \\
\\ 
\texttt{If the user takes the bucket and battery (2 items), then they can take one knife, making 3 items. The other player would then have 2 knives left. Alternatively, if the user takes all three knives, that's 3 items, leaving the bucket and battery for the other player. But if the other player wants both the bucket and battery, that's two items, which would be complementary. However, if the other player also wants some knives, taking all three might not be complementary.} \\
            }
        }
    }
    & & \\ \\

    \theutterance \stepcounter{utterance}  
    & & \multicolumn{4}{p{0.6\linewidth}}{
        \cellcolor[rgb]{0.9,0.9,0.9}{
            \makecell[{{p{\linewidth}}}]{
                \texttt{\tiny{[GM$\rangle$P2]}}
                \texttt{<think>} \\
\texttt{Okay, let's try to figure this out. The user is trying to divide the items between themselves and another player. The items are 1 bucket, 3 knives, and 1 battery. The user's value function is 2 for each item, so they value all items equally. The other player's values are secret, so the user doesn't know what they prefer.} \\
\\ 
\texttt{The user's goal is to maximize their own score, which is the sum of the values of the items they receive. If their proposal and the other player's proposal are complementary (i.e., they don't overlap and all items are distributed), then both get their scores. Otherwise, both get zero.} \\
\\ 
\texttt{In the previous attempt, the user considered taking the bucket, battery, and one knife, leaving two knives for the other player. That would give the user 3 items (3*2=6 points) and the other player 2 items (2*2=4 points). But the user is unsure if this is the best approach because they don't know the other player's preferences.} \\
\\ 
\texttt{Another possibility is to take more knives. Since there are three knives, maybe taking two knives and one other item. For example, two knives and the bucket. That would leave one knife and the battery for the other player. But again, the user doesn't know if the other player values the battery or bucket more.} \\
\\ 
\texttt{Wait, the user's value is the same for all items, so they just want as many items as possible. The maximum number of items they can get without overlapping with the other player's proposal is 3 (since there are 5 items total). So if the user takes 3 items, the other player can take 2. But the user needs to make sure that the other player's proposal doesn't conflict. Since the other player's values are unknown, the user has to make a proposal that is likely to be complementary regardless of the other player's preferences.} \\
\\ 
\texttt{If the user takes the bucket and battery (2 items), then they can take one knife, making 3 items. The other player would then have 2 knives left. Alternatively, if the user takes all three knives, that's 3 items, leaving the bucket and battery for the other player. But if the other player wants both the bucket and battery, that's two items, which would be complementary. However, if the other player also wants some knives, taking all three might not be complementary.} \\
            }
        }
    }
    & \\ \\

    \theutterance \stepcounter{utterance}  
    & & & \multicolumn{4}{p{0.6\linewidth}}{
        \cellcolor[rgb]{0.9,0.9,0.9}{
            \makecell[{{p{\linewidth}}}]{
                \texttt{\tiny{[GM$\langle$P2]}}
                \texttt{<think>} \\
\texttt{Okay, let's try to figure this out. The user is trying to divide the items between themselves and another player. The items are 1 bucket, 3 knives, and 1 battery. The user's value function is 2 for each item, so they value all items equally. The other player's values are secret, so the user doesn't know what they prefer.} \\
\\ 
\texttt{The user's goal is to maximize their own score, which is the sum of the values of the items they receive. If their proposal and the other player's proposal are complementary (i.e., they don't overlap and all items are distributed), then both get their scores. Otherwise, both get zero.} \\
\\ 
\texttt{In the previous attempt, the user considered taking the bucket, battery, and one knife, leaving two knives for the other player. That would give the user 3 items (3*2=6 points) and the other player 2 items (2*2=4 points). But the user is unsure if this is the best approach because they don't know the other player's preferences.} \\
\\ 
\texttt{Another possibility is to take more knives. Since there are three knives, maybe taking two knives and one other item. For example, two knives and the bucket. That would leave one knife and the battery for the other player. But again, the user doesn't know if the other player values the battery or bucket more.} \\
\\ 
\texttt{Wait, the user's value is the same for all items, so they just want as many items as possible. The maximum number of items they can get without overlapping with the other player's proposal is 3 (since there are 5 items total). So if the user takes 3 items, the other player can take 2. But the user needs to make sure that the other player's proposal doesn't conflict. Since the other player's values are unknown, the user has to make a proposal that is likely to be complementary regardless of the other player's preferences.} \\
\\ 
\texttt{If the user takes the bucket and battery (2 items), then they can take one knife, making 3 items. The other player would then have 2 knives left. Alternatively, if the user takes all three knives, that's 3 items, leaving the bucket and battery for the other player. But if the other player wants both the bucket and battery, that's two items, which would be complementary. However, if the other player also wants some knives, taking all three might not be complementary.} \\
            }
        }
    }
     \\ \\

    \theutterance \stepcounter{utterance}  
    & & & \multicolumn{2}{p{0.3\linewidth}}{
        \cellcolor[rgb]{0.9,0.9,0.9}{
            \makecell[{{p{\linewidth}}}]{
                \texttt{\tiny{[GM$|$GM]}}
                \texttt{continue} \\
            }
        }
    }
    & & \\ \\

    \theutterance \stepcounter{utterance}  
    & & & \multicolumn{2}{p{0.3\linewidth}}{
        \cellcolor[rgb]{0.9,0.9,0.9}{
            \makecell[{{p{\linewidth}}}]{
                \texttt{\tiny{[GM$|$GM]}}
                \texttt{<think>} \\
\texttt{Okay, let's try to figure this out. The user is trying to divide the items between themselves and another player. The items are 1 bucket, 3 knives, and 1 battery. The user's value function is 2 for each item, so they value all items equally. The other player's values are secret, so the user doesn't know what they prefer.} \\
\\ 
\texttt{The user's goal is to maximize their own score, which is the sum of the values of the items they receive. If their proposal and the other player's proposal are complementary (i.e., they don't overlap and all items are distributed), then both get their scores. Otherwise, both get zero.} \\
\\ 
\texttt{In the previous attempt, the user considered taking the bucket, battery, and one knife, leaving two knives for the other player. That would give the user 3 items (3*2=6 points) and the other player 2 items (2*2=4 points). But the user is unsure if this is the best approach because they don't know the other player's preferences.} \\
\\ 
\texttt{Another possibility is to take more knives. Since there are three knives, maybe taking two knives and one other item. For example, two knives and the bucket. That would leave one knife and the battery for the other player. But again, the user doesn't know if the other player values the battery or bucket more.} \\
\\ 
\texttt{Wait, the user's value is the same for all items, so they just want as many items as possible. The maximum number of items they can get without overlapping with the other player's proposal is 3 (since there are 5 items total). So if the user takes 3 items, the other player can take 2. But the user needs to make sure that the other player's proposal doesn't conflict. Since the other player's values are unknown, the user has to make a proposal that is likely to be complementary regardless of the other player's preferences.} \\
\\ 
\texttt{If the user takes the bucket and battery (2 items), then they can take one knife, making 3 items. The other player would then have 2 knives left. Alternatively, if the user takes all three knives, that's 3 items, leaving the bucket and battery for the other player. But if the other player wants both the bucket and battery, that's two items, which would be complementary. However, if the other player also wants some knives, taking all three might not be complementary.} \\
            }
        }
    }
    & & \\ \\

    \theutterance \stepcounter{utterance}  
    & & \multicolumn{4}{p{0.6\linewidth}}{
        \cellcolor[rgb]{0.9,0.9,0.9}{
            \makecell[{{p{\linewidth}}}]{
                \texttt{\tiny{[P1$\langle$GM]}}
                \texttt{<think>} \\
\texttt{Okay, let's try to figure this out. The user is trying to divide the items between themselves and another player. The items are 1 bucket, 3 knives, and 1 battery. The user's value function is 2 for each item, so they value all items equally. The other player's values are secret, so the user doesn't know what they prefer.} \\
\\ 
\texttt{The user's goal is to maximize their own score, which is the sum of the values of the items they receive. If their proposal and the other player's proposal are complementary (i.e., they don't overlap and all items are distributed), then both get their scores. Otherwise, both get zero.} \\
\\ 
\texttt{In the previous attempt, the user considered taking the bucket, battery, and one knife, leaving two knives for the other player. That would give the user 3 items (3*2=6 points) and the other player 2 items (2*2=4 points). But the user is unsure if this is the best approach because they don't know the other player's preferences.} \\
\\ 
\texttt{Another possibility is to take more knives. Since there are three knives, maybe taking two knives and one other item. For example, two knives and the bucket. That would leave one knife and the battery for the other player. But again, the user doesn't know if the other player values the battery or bucket more.} \\
\\ 
\texttt{Wait, the user's value is the same for all items, so they just want as many items as possible. The maximum number of items they can get without overlapping with the other player's proposal is 3 (since there are 5 items total). So if the user takes 3 items, the other player can take 2. But the user needs to make sure that the other player's proposal doesn't conflict. Since the other player's values are unknown, the user has to make a proposal that is likely to be complementary regardless of the other player's preferences.} \\
\\ 
\texttt{If the user takes the bucket and battery (2 items), then they can take one knife, making 3 items. The other player would then have 2 knives left. Alternatively, if the user takes all three knives, that's 3 items, leaving the bucket and battery for the other player. But if the other player wants both the bucket and battery, that's two items, which would be complementary. However, if the other player also wants some knives, taking all three might not be complementary.} \\
\\ 
\\ 
\texttt{The time is up. It is now your turn to submit a secret proposal.} \\
            }
        }
    }
    & \\ \\

    \theutterance \stepcounter{utterance}  
    & \multicolumn{4}{p{0.6\linewidth}}{
        \cellcolor[rgb]{0.9,0.9,0.9}{
            \makecell[{{p{\linewidth}}}]{
                \texttt{\tiny{[P1$\rangle$GM]}}
                \texttt{<think>} \\
\texttt{Okay, let's try to figure this out. The user is trying to divide the items between themselves and another player. The items are 1 bucket, 3 knives, and 1 battery. The user's value function is 2 for each item, so they value all items equally. The other player's values are secret, so the user doesn't know what they prefer.} \\
\\ 
\texttt{The user's goal is to maximize their own score, which is the sum of the values of the items they receive. If their proposal and the other player's proposal are complementary (i.e., they don't overlap and all items are distributed), then both get their scores. Otherwise, both get zero.} \\
\\ 
\texttt{In the previous attempt, the user considered taking the bucket, battery, and one knife, leaving two knives for the other player. That would give the user 3 items (3*2=6 points) and the other player 2 items (2*2=4 points). But the user is unsure if this is the best approach because they don't know the other player's preferences.} \\
\\ 
\texttt{Another possibility is to take more knives. Since there are three knives, maybe taking two knives and one other item. For example, two knives and the bucket. That would leave one knife and the battery for the other player. But again, the user doesn't know if the other player values the battery or bucket more.} \\
\\ 
\texttt{Wait, the user's value is the same for all items, so they just want as many items as possible. The maximum number of items they can get without overlapping with the other player's proposal is 3 (since there are 5 items total). So if the user takes 3 items, the other player can take 2. But the user needs to make sure that the other player's proposal doesn't conflict. Since the other player's values are unknown, the user has to make a proposal that is likely to be complementary regardless of the other player's preferences.} \\
\\ 
\texttt{If the user takes the bucket and battery (2 items), then they can take one knife, making 3 items. The other player would then have 2 knives left. Alternatively, if the user takes all three knives, that's 3 items, leaving the bucket and battery for the other player. But if the other player wants both the bucket and battery, that's two items, which would be complementary. However, if the other player also wants some knives, taking all three might not be complementary.} \\
            }
        }
    }
    & & \\ \\

    \theutterance \stepcounter{utterance}  
    & & & \multicolumn{2}{p{0.3\linewidth}}{
        \cellcolor[rgb]{0.9,0.9,0.9}{
            \makecell[{{p{\linewidth}}}]{
                \texttt{\tiny{[GM$|$GM]}}
                \texttt{continue} \\
            }
        }
    }
    & & \\ \\

    \theutterance \stepcounter{utterance}  
    & & & \multicolumn{2}{p{0.3\linewidth}}{
        \cellcolor[rgb]{0.9,0.9,0.9}{
            \makecell[{{p{\linewidth}}}]{
                \texttt{\tiny{[GM$|$GM]}}
                \texttt{player was instructed to make a proposal} \\
            }
        }
    }
    & & \\ \\

\end{supertabular}
}

\end{document}
